\section{Replace Elements by Rank in the Array}
\label{sec:heap-replace-by-rank}

\problemstatement{Given an array $\textit{arr}$ of $n$ integers, replace
every element with its \textbf{rank} when the array's distinct values are
sorted in increasing order: the smallest value gets rank $1$, the next
distinct value gets rank $2$, and so on; equal elements receive the same
rank.}

\begin{patternbox}
\emph{``Replace each element by its rank in sorted order''} is, at its
core, a sorting problem -- but presenting it via a min-heap makes the
\textbf{``process elements in increasing order while remembering where
each one came from''} pattern explicit, which is the same underlying idea
used throughout this chapter whenever output order must differ from input
order. Whenever you need to process values in sorted order \emph{and}
write results back to their original positions, pair each value with its
original index before pushing it into the heap.
\end{patternbox}

\subsection*{Understanding the problem carefully}

Rank assignment requires processing values from smallest to largest while
remembering, for each value, which original array position it came from
(so the computed rank can be written back to the correct place). A
min-heap of $(\textit{value}, \textit{originalIndex})$ pairs, ordered by
value, gives exactly this controlled smallest-to-largest traversal.

\begin{ideabox}[Greedy Choice, Powered by a Min-Heap]
Push every $(\textit{value}, \textit{originalIndex})$ pair into a
min-heap ordered by value. Maintain a running rank counter, starting at
$0$, and a record of the previously popped value (initially undefined).
Repeatedly pop the smallest remaining pair: if its value differs from the
previously popped value (or this is the first pop), increment the rank
counter; write the current rank counter's value into the result array at
the popped pair's original index.
\end{ideabox}

\subsection*{Why incrementing the rank only on a value change is correct}

Equal elements must receive the same rank, and the rank must strictly
increase only when a genuinely new, larger distinct value is encountered.
Since the heap pops values in non-decreasing order, comparing each popped
value to the immediately previous one correctly detects exactly the
moments where a new distinct value begins -- there is no need to
pre-deduplicate the array, because consecutive equal pops are detected and
handled identically (no rank increment) regardless of how many times a
value repeats.

\subsection*{Algorithm}

\begin{enumerate}
  \item Push every $(\textit{arr}[i], i)$ pair onto a min-heap, ordered
        by value.
  \item Initialize $\textit{rank} \gets 0$, $\textit{prevValue} \gets$
        undefined.
  \item While the heap is not empty:
  \begin{enumerate}
    \item Pop $(\textit{value}, \textit{idx})$.
    \item If $\textit{value} \ne \textit{prevValue}$ (or this is the
          first pop): $\textit{rank} \gets \textit{rank}+1$;
          $\textit{prevValue} \gets \textit{value}$.
    \item $\textit{result}[\textit{idx}] \gets \textit{rank}$.
  \end{enumerate}
  \item Return \textit{result}.
\end{enumerate}

\begin{algorithm}[H]
\caption{Replace Elements by Rank}
\begin{algorithmic}[1]
\Procedure{ReplaceByRank}{\textit{arr}}
  \State $\textit{heap} \gets$ min-heap of $(\textit{arr}[i], i)$ for all $i$
  \State $\textit{rank} \gets 0$, $\textit{prevValue} \gets$ undefined
  \While{\textit{heap} is not empty}
    \State $(\textit{value}, \textit{idx}) \gets$ pop minimum
    \If{$\textit{value} \ne \textit{prevValue}$}
      \State $\textit{rank} \gets \textit{rank} + 1$
      \State $\textit{prevValue} \gets \textit{value}$
    \EndIf
    \State $\textit{result}[\textit{idx}] \gets \textit{rank}$
  \EndWhile
  \State \Return \textit{result}
\EndProcedure
\end{algorithmic}
\end{algorithm}

\subsection*{Dry run}

\dryrun{Take $\textit{arr} = [20, 15, 26, 2, 98, 6]$.}

Heap pairs: $(20,0),(15,1),(26,2),(2,3),(98,4),(6,5)$.

\begin{itemize}
  \item Pop $(2,3)$: new value. $\textit{rank}=1$.
        $\textit{result}[3]=1$.
  \item Pop $(6,5)$: new value. $\textit{rank}=2$.
        $\textit{result}[5]=2$.
  \item Pop $(15,1)$: new value. $\textit{rank}=3$.
        $\textit{result}[1]=3$.
  \item Pop $(20,0)$: new value. $\textit{rank}=4$.
        $\textit{result}[0]=4$.
  \item Pop $(26,2)$: new value. $\textit{rank}=5$.
        $\textit{result}[2]=5$.
  \item Pop $(98,4)$: new value. $\textit{rank}=6$.
        $\textit{result}[4]=6$.
\end{itemize}

Final result: $[4,3,5,1,6,2]$.

\subsection*{C++ implementation}

\begin{lstlisting}
#include <bits/stdc++.h>
using namespace std;

vector<int> replaceByRank(vector<int>& arr) {
    int n = arr.size();
    priority_queue<pair<int,int>, vector<pair<int,int>>,
                   greater<pair<int,int>>> minHeap;

    for (int i = 0; i < n; i++) {
        minHeap.push({arr[i], i});
    }

    vector<int> result(n);
    int rank = 0;
    bool first = true;
    int prevValue = -1;

    while (!minHeap.empty()) {
        auto [value, idx] = minHeap.top();
        minHeap.pop();

        if (first || value != prevValue) {
            rank++;
            prevValue = value;
            first = false;
        }
        result[idx] = rank;
    }
    return result;
}

int main() {
    vector<int> arr = {20, 15, 26, 2, 98, 6};
    vector<int> result = replaceByRank(arr);

    for (int x : result) cout << x << " ";
    cout << endl;
    return 0;
}
\end{lstlisting}

\begin{complexitybox}
\textbf{Time Complexity.} Pushing $n$ pairs costs $O(n \log n)$, and
popping all $n$ pairs costs another $O(n \log n)$, giving an overall time
complexity of
\[
  O(n \log n).
\]
\textbf{Space Complexity.} $O(n)$ for the heap and the result array.
\end{complexitybox}

\begin{takeawaybox}
Whenever a problem needs values processed in sorted order while results
must land back at their \emph{original} positions, push
$(\textit{value}, \textit{originalIndex})$ pairs together, rather than
sorting the values alone and losing track of where each one belongs. This
small habit -- carry the index along with the value -- recurs constantly
whenever sorting (or heap order) and output position diverge.
\end{takeawaybox}
